\chapter{Conclusion}\label{chap:concl}

This thesis developed a single closed-loop adaptive tutoring system in which \ac{FER} and \ac{RL} function as inseparable layers of one control architecture rather than as adjacent models. A commodity webcam supplies the perception layer; a multi-label academic \ac{FER} component trained on the \ac{DAiSEE} dataset projects the learner's affective state onto a continuous signal; and a \ac{DQN} decision layer consumes that signal directly to select pedagogical actions within a theory-grounded Markov decision process over an eight-action space. The two layers are coupled through the \texttt{FERAdapter}, a unifying abstraction that lets the same policy operate over either simulated or real affect without retraining. The central result is direct: an affect-aware policy raised normalised learning gain by $0.088$ over a hand-crafted, affect-keyed expert heuristic on a heterogeneous population of synthetic learners (Cohen's $d = 2.01$, Holm-corrected $p < 10^{-2}$), and held that advantage across every outcome metric in both populations, with effect sizes from $1.81$ to $3.30$. Intelligent tutoring systems have modelled what learners know with increasing sophistication while ignoring how they feel; these simulation results suggest that incorporating affect into instructional decisions can improve modelled learning outcomes within the synthetic tutoring environment used in this thesis. Whether these improvements transfer to real learners remains an open question for future work.

The perception layer earns that claim by first repairing the ground beneath it. Published \ac{DAiSEE} results are mutually incomparable, because preprocessing, frame counts, and balancing differ silently between papers. Retraining six competing architectures from scratch under one fixed pipeline turned that literature into a controlled comparison, and within it the proposed ResNet-50 with Squeeze-and-Excitation recalibration, a two-layer \ac{BiLSTM}, and temporal attention pooling won, reaching the highest Macro F1 at $0.453$ and a mean per-label accuracy of $0.782$. Framing academic affect as four independent binary labels rather than one dominant class let the model represent the engaged-yet-confused learner that a single-label head cannot, and the per-label F1 scores record both the achievement and its honest boundary: $0.975$ on engagement, $0.358$ on boredom, $0.255$ on frustration, and $0.224$ on confusion. The decision layer earns its claim through restraint. Every transition rule in the environment traces to published theory, and the \ac{DQN} is the only component allowed to learn, so the policy's learning gain of $0.548$ against $0.460$ for the expert heuristic reflects learned control, not a simulator bent in its favour. The generalisation result carries the most weight: a policy trained on diverse learners transferred upward to the canonical learner at $0.607$, while a policy trained on the canonical learner alone fell to $0.553$ on the mixed population in every one of ten seeds. Robustness in affect-aware tutoring comes from exposure to learner variability, not from architectural depth, and the system that demonstrates it remains lightweight enough to deploy on commodity hardware.

This thesis shows that benchmark-validated affect perception and theory-grounded decision-making can be joined through a single interface into one reproducible architecture, with each half measured honestly and the connection between them made explicit. For intelligent tutoring, it shows that the affective signal a good teacher reads by instinct can be instrumented, learned, and acted on without surrendering rigour. For affective computing, it shows that emotion recognition pays off not as a dashboard but as a control signal inside a closed loop. For personalised education at scale, it shows that tutors worthy of real learners are built by training on the full variability of the people they will serve. The learner's face has always been the richest signal in the room; this work demonstrates in simulation, in numbers and in architecture, that a machine can be built to read it, reason from it, and earn a place beside the teacher rather than the test.